% Methodology chapter. Settings are those in configs/*.yaml; the implementation is pinn_tailrisk/.
\chapter{Methodology}
\label{ch:methodology}

This chapter describes the model (\cref{sec:met-model}), the reference methods that provide ground truth
(\cref{sec:met-reference}), the physics-informed network in one and in $d$ dimensions
(\cref{sec:met-pinn1d,sec:met-pinnnd}), how risk is read from it (\cref{sec:met-risk}), the amortised
model (\cref{sec:met-amortised}), the real-data study (\cref{sec:met-real}), the data-only baseline
(\cref{sec:met-dataonly}) and the evaluation protocol (\cref{sec:met-eval}). Time is measured in years of
252 trading days and returns are log returns throughout.

% ------------------------------------------------------------------------------------------------
\section{Model}
\label{sec:met-model}

\subsection{One asset}
The log return $X_t$ (with $X_0 = 0$) follows the Merton jump-diffusion
\begin{equation}
  dX_t = a\,dt + \sigma\,dW_t + dJ_t, \qquad J_t = \sum_{i=1}^{N_t} Y_i, \qquad
  N_t \sim \mathrm{Poisson}(\lambda t), \quad Y_i \sim \mathcal{N}(\mu_J, \sigma_J^2).
  \label{eq:merton}
\end{equation}
Its density $p(x,t)$ solves the forward Kolmogorov (Fokker--Planck) partial integro-differential equation
\begin{equation}
  p_t = -a\,p_x + \tfrac12 \sigma^2 p_{xx} + \lambda \int_{\mathbb{R}} p(x - y, t)\, f_Y(y)\,dy - \lambda\, p,
  \label{eq:pide1d}
\end{equation}
with $p(\cdot, 0) = \delta_0$. The exact solution is a Poisson mixture of Gaussians \citep{merton1976},
\begin{equation}
  p(x,t) = \sum_{k \ge 0} \frac{e^{-\lambda t}(\lambda t)^k}{k!}\,
           \mathcal{N}\!\left(x;\; a t + k\mu_J,\; \sigma^2 t + k\sigma_J^2\right),
  \label{eq:exact1d}
\end{equation}
which is what makes the model suitable for a validation study: every PINN output can be compared with an
exact answer. The base parameters, representing an equity index with a crash component, are $a = 0.05$,
$\sigma = 0.20$, $\lambda = 5$/yr, $\mu_J = -0.05$ and $\sigma_J = 0.07$.

\subsection{Several assets with a common jump}
For $d$ assets the diffusion is $\mathcal{N}(a t, \Sigma t)$ and all assets share one market-jump factor:
on a jump every asset moves by $\beta_i Y$ with a scalar $Y \sim \mathcal{N}(\mu_J, \sigma_J^2)$ and fixed
loadings $\beta \in \mathbb{R}^d$. The forward equation becomes
\begin{equation}
  p_t = -a \cdot \nabla p + \tfrac12 \operatorname{tr}\!\big(\Sigma \nabla^2 p\big)
        + \lambda \int_{\mathbb{R}} p(x - \beta y, t)\, f_Y(y)\,dy - \lambda\, p .
  \label{eq:pidend}
\end{equation}
Two properties follow from the common-jump structure. First, the jump integral is one-dimensional in any
dimension, so it can be evaluated by the same quadrature as in 1-D; a full $d$-dimensional jump distribution
would make it a $d$-dimensional integral. Second, the density is again an exact Gaussian mixture, with
component $k$ having mean $a t + k\mu_J\beta$ and covariance $\Sigma t + k\sigma_J^2\beta\beta^\top$.

The family used in the scaling study has volatilities spread linearly over $[0.18, 0.30]$, loadings over
$[0.8, 1.2]$ and equicorrelation $\rho = 0.4$, with the jump parameters of the 1-D model.

\subsection{Why the multi-asset density is needed}
For a linear portfolio with weights $w$, the portfolio log return $w^\top X$ is itself a 1-D Merton process
with drift $w^\top a$, volatility $\sqrt{w^\top \Sigma w}$ and jump $(w^\top \beta)\,Y$. A linear portfolio
alone therefore does not require a $d$-dimensional solver. We use this projection as an exact reference for
linear portfolios, and the joint density for what the projection cannot provide: one density that serves every
portfolio and gives risk contributions by differentiation (\cref{sec:met-amortised}), and non-linear
portfolios that hold options (\cref{sec:met-real}).

% ------------------------------------------------------------------------------------------------
\section{Reference methods}
\label{sec:met-reference}

\paragraph{Closed form.} \Cref{eq:exact1d} and its multivariate version, truncated where the neglected
Poisson mass is negligible, give exact densities. Exact VaR and ES of a linear portfolio come from the 1-D
projection by root finding on the mixture CDF and the closed-form partial expectation of each Gaussian
component.

\paragraph{Finite differences in 1-D.} \Cref{eq:pide1d} is solved by an IMEX scheme: Crank--Nicolson for
drift, diffusion and $-\lambda p$, and second-order Adams--Bashforth for the jump convolution, which is applied
by FFT with a cell-averaged jump kernel so that mass is exact on the grid. The scheme is second order in space
and time. It runs on $[-1.0, 0.5]$ with 4001 points and $p = 0$ at both ends; this domain is wider than the one
the PINN is evaluated on, because on a narrower one the zero boundaries distort the far tail.

\paragraph{Finite differences in 2-D.} For $d = 2$ a Douglas ADI scheme ($\theta = 1/2$) treats the
axis-aligned drift and diffusion implicitly with one tridiagonal sweep per axis, and the mixed derivative and
the jump integral explicitly. The jump kernel lies on the line $\{\beta y\}$; it is sampled finely, spread
bilinearly onto grid offsets (preserving mass) and applied by FFT, so a time step costs $O(n^2 \log n)$. This
solver is a second, independent reference in 2-D and a measurement of how grid cost grows with dimension.

\paragraph{Monte Carlo.} Exact simulation of \cref{eq:merton}: a Poisson number of jumps, then the
Gaussian conditional on it. VaR and ES are the empirical quantile and tail mean.

% ------------------------------------------------------------------------------------------------
\section{The PINN in one dimension}
\label{sec:met-pinn1d}

\paragraph{Output and inputs.} The network $g_\phi$ (four hidden layers of width 64, $\tanh$) returns the
\emph{log} density,
\begin{equation}
  u(x,t) = \log p(x,t) = -\tfrac12 \log(2\pi\sigma^2 t) + c\, g_\phi(z, \tau), \qquad
  z = \frac{x - a t}{\sigma\sqrt{t}}, \qquad
  \tau = 2\,\frac{\log t - \log t_0}{\log T - \log t_0} - 1,
  \label{eq:u1d}
\end{equation}
with output scale $c = 10$. Three reasons motivate this form. The density is positive by construction.
Tail densities of order $10^{-4}$ become numbers of order one, which a network can fit. And in the similarity
variable $z$ the diffusive bulk has unit width at every time, instead of collapsing as $t \to t_0$; the known
Gaussian normaliser is added outside the network.

\paragraph{Log-form residual.} Dividing \cref{eq:pide1d} by $p$ gives an equation for $u$,
\begin{equation}
  r(x,t) = u_t + a\,u_x - \tfrac12\sigma^2\big(u_{xx} + u_x^2\big) - \lambda\,(S - 1), \qquad
  S(x,t) = \int e^{\,u(x-y,t) - u(x,t)}\, f_Y(y)\,dy .
  \label{eq:logres}
\end{equation}
The standard residual (on $p$) weighs a point by the size of the density there, so the tail, where risk is
determined, contributes almost nothing to the loss. The log-form residual measures relative error and weighs a
tail point as much as a bulk point. The residual is multiplied by $t$ to make it dimensionless across the time
range.

\paragraph{Jump integral.} Since $Y$ is Gaussian, $S$ is computed by Gauss--Hermite quadrature with nodes
$y_j = \mu_J + \sqrt{2}\,\sigma_J \xi_j$ and weights $w_j/\sqrt{\pi}$; 48 nodes are used (chosen by the ablation
in the results chapter).

\paragraph{Initial condition.} Training starts at $t_0 = 0.5$ day from the exact density
\cref{eq:exact1d}, which avoids the Dirac delta at $t = 0$. In practice this short-time density would come from
the diffusion approximation plus at most one jump; here it is exact so that the study measures the PDE solve
alone.

\paragraph{Loss.} The loss is
\begin{equation}
  \mathcal{L} = w_{\mathrm{pde}}\,\overline{r^2}
              + w_{\mathrm{ic}}\,\overline{\big(u(x, t_0) - \log p(x, t_0)\big)^2}
              + w_{\mathrm{norm}}\,\overline{\Big(\log \textstyle\int e^{u(x,t)}\,dx\Big)^2}
              + w_{\mathrm{bc}}\,\mathcal{L}_{\mathrm{decay}},
  \label{eq:loss1d}
\end{equation}
with weights 1, 10, 10 and 1. The normalisation term keeps $\int p = 1$, which the residual of
\cref{eq:logres} does not enforce on its own (it is unchanged by adding a constant to $u$). The decay term
asks the density to fall at the ends of the domain $[-0.6, 0.3]$. In 1-D, 4096 collocation points are drawn
per step, 40\% of them around the model's current 1\% quantile (tail-adaptive sampling, refreshed every 250
steps).

\paragraph{Optimisation.} Adam with learning rate $2\times10^{-3}$, cosine-annealed to $2\times10^{-5}$ over
20{,}000 steps, with gradient clipping at norm 1.

% ------------------------------------------------------------------------------------------------
\section{The PINN in $d$ dimensions}
\label{sec:met-pinnnd}

The 1-D design carries over with four changes.

\paragraph{Whitened input.} With $L$ the Cholesky factor of $\Sigma$, the network sees
$z = L^{-1}(x - a t)/\sqrt{t}$ and $\tau$, and the Gaussian normaliser
$-\tfrac12\log\det(2\pi\Sigma t)$ is added to its output. The network has width 128 and depth 4.

\paragraph{Residual.} The log-form residual of \cref{eq:pidend} is
\begin{equation}
  r = u_t + a\cdot\nabla u - \tfrac12\Big[\operatorname{tr}\!\big(\Sigma\nabla^2 u\big)
      + \nabla u^\top \Sigma\, \nabla u\Big] - \lambda (S - 1), \qquad
  S = \int e^{\,u(x - \beta y) - u(x)} f_Y(y)\,dy,
\end{equation}
and $S$ uses the same 48-node quadrature along $\beta$.

\paragraph{Trace term.} $\operatorname{tr}(\Sigma\nabla^2 u) = \sum_i \ell_i^\top \nabla^2 u\, \ell_i$, with
$\ell_i$ the columns of $L$, is computed exactly by $d$ Hessian-vector products for $d \le 4$. For larger $d$ it
is replaced by Hutchinson's estimator $\mathbb{E}_v\big[(Lv)^\top \nabla^2 u\,(Lv)\big]$ with four Rademacher
probes $v$, whose cost does not grow with $d$.

\paragraph{No grid: a proposal built from the coefficients.} A box domain cannot be used in $d$ dimensions,
so collocation, the normalisation integral and the risk computation all draw from a proposal
\begin{equation}
  q(x \mid t) = \sum_{k=0}^{K} \pi_k(t)\; t_\nu\!\Big(x;\; a t + k\mu_J\beta,\;
                c^2\big(\Sigma t + k\sigma_J^2 \beta\beta^\top\big)\Big),
  \label{eq:proposal}
\end{equation}
a mixture over $k = 0,\dots,K$ jumps of multivariate Student-$t$ components. The weights $\pi_k$ are the Poisson
probabilities, floored at 0.05 and renormalised so that every jump count is sampled. The proposal uses only the
coefficients of the PIDE, not its solution. Two versions are used: for collocation and the initial condition,
$\nu = 30$ and $c = 1.5$ (heavier tails put training targets at $\log p \approx -7000$ in high dimension); for
the normalisation integral and for risk, $\nu = 5$ and $c = 1.2$, so that importance weights stay bounded where
the density is non-negligible. The default is $K = 3$ jumps. Tail-adaptive sampling is not used in $d$
dimensions, since it gave no gain in 1-D.

\paragraph{Normalisation loss.} With $M$ proposal draws at each of several times,
\begin{equation}
  \mathcal{L}_{\mathrm{norm}} = \overline{\Big(\log \tfrac1M \textstyle\sum_m e^{\,u(x_m,t) - \log q(x_m \mid t)}\Big)^2},
\end{equation}
computed with a log-sum-exp. Penalising $(\text{mass} - 1)^2$ directly overflows for an untrained network.
The loss weights, optimiser and schedule are those of the 1-D model, with 2048 collocation points, 4096
initial-condition points and 4096 normalisation draws per step.

% ------------------------------------------------------------------------------------------------
\section{Risk from the network}
\label{sec:met-risk}

The loss is $L = -X$ and, at confidence level $\alpha$,
$\VaR_\alpha = \inf\{\ell : P(L \le \ell) \ge \alpha\}$ and $\ES_\alpha = \mathbb{E}[L \mid L \ge \VaR_\alpha]$,
both reported as positive numbers. In 1-D they are read from the density on a grid. In $d$ dimensions they are
computed by self-normalised importance sampling: draws $x_m \sim q(\cdot \mid t)$ from \cref{eq:proposal} get
weights $\tilde w_m \propto \exp\big(u(x_m,t) - \log q(x_m \mid t)\big)$, and VaR and ES of any portfolio
$P\&L(x)$ are the weighted quantile and the weighted tail mean of $P\&L(x_m)$, with the atom at the quantile
split exactly. The same draws serve every portfolio. Because the density is an explicit function of $x$, $w$
and (in the amortised model) the market parameters, risk contributions $\partial\ES/\partial w_i$ follow from
the weighted tail average of each asset's return.

\paragraph{Trust region.} The network is constrained only where collocation points are drawn. Far outside that
region it can decay more slowly than the proposal, and a single draw then dominates the weights. For the
real-data study, draws whose Mahalanobis radius to the nearest collocation component exceeds $r_{\max} = 6$ are
given zero weight; the exact mass removed in this way is below $10^{-15}$.

% ------------------------------------------------------------------------------------------------
\section{Amortised model}
\label{sec:met-amortised}

To answer queries for a changing market without retraining, the network takes the parameters
$\theta = (s, \lambda, \mu_J)$ as additional inputs. Here $s$ multiplies every diffusion volatility
(covariance $s^2\Sigma$, with correlations and loadings fixed), $\lambda$ is the jump intensity and $\mu_J$
the mean jump; $\sigma_J$ is fixed. The whitened input is also divided by $s$. Training draws a fresh $\theta$
for every collocation, initial-condition and normalisation point, from the box $s \in [0.6, 1.6]$,
$\lambda \in [1, 15]$/yr (log-uniform) and $\mu_J \in [-0.10, -0.02]$. The initial condition is the exact mixture
at $t_0$ evaluated at each point's $\theta$, and the proposal of \cref{eq:proposal} is evaluated at each row's
$\theta$. The exact trace is used, and training runs for 30{,}000 steps.

Latency is measured on the same device for PINN importance sampling with $N_{\mathrm{IS}}$ draws and Monte
Carlo with $N_{\mathrm{MC}}$ paths. Fitting error $\propto N^{\text{slope}}$ and latency $\approx l_0 + l_1 N$
for each method gives the cost of each at the PINN's error level, and from that the number of queries at which
the training cost is paid back (the break-even).

% ------------------------------------------------------------------------------------------------
\section{Real-data study}
\label{sec:met-real}

\paragraph{Data.} Two five-asset baskets: the five largest NIFTY 50 constituents (Reliance, HDFC Bank, Infosys,
ICICI Bank, TCS) and five S\&P 500 sector ETFs (XLK, XLF, XLE, XLV, XLI). Daily adjusted prices from Yahoo
Finance from January 2006 (earlier NIFTY prices contain bad prints) to October 2026. The log returns are
committed with the code, since the provider revises adjusted prices. Models are developed on 2006--2017 and
tested from January 2018, which includes the COVID crash.

\paragraph{Calibration.} The common-jump model is fitted by maximum likelihood to daily returns, each day being
one draw of $X$ over one trading day with density the Poisson mixture over $k \le 8$ jumps. In the
\emph{filtered} model a daily volatility scale $s_t$ from an EWMA filter (decay 0.94) multiplies every diffusion
volatility, which is exactly the parameter the amortised PINN takes as input. Jump loadings are fixed equal: with
free loadings the NIFTY jump attaches to a single stock and stops being a market jump. Parameters are optimised
jointly with L-BFGS in double precision, and standard errors come from the inverse Hessian. During the test
period $(\lambda, \mu_J)$ are refitted monthly on a 1000-day window with a Gaussian prior centred on the
development fit (standard deviation 0.5 for $\log\lambda$ and 0.01 for $\mu_J$), because the rolling fit alone is
weakly identified. An \emph{unfiltered} model (no $s_t$, all parameters refitted monthly) is a separate
benchmark.

\paragraph{The PINN.} One amortised network per market is trained on a box set from development data only:
$s$ over the range seen in development (widened by factors 0.75 and 1.35), $\lambda$ within a factor $e^{\pm1.5}$ of
the development fit and $\mu_J$ within $\pm 0.025$. Time runs over $[0.5, 1.5]$ days, since only one-day-ahead
forecasts are needed, and $K = 8$ jumps are allowed because real-data jumps are small and frequent. Each test day
is one forward query at that day's $\theta_t$, using 32{,}768 importance draws. The same model evaluated with its
exact density on the same draws is reported alongside, so the difference between the two is the PINN's error
and not the model's.

\paragraph{Portfolios.} An equal-weight basket (linear) and the basket plus a short one-month put struck 3\%
out of the money, rolled monthly (non-linear). The put is priced by Black--Scholes with zero rate and implied
volatility equal to the EWMA volatility of the basket, and the book is fully revalued on every scenario.

\paragraph{Benchmarks.} All forecast one day ahead without look-ahead: Gaussian (250-day window), historical
simulation (500 days), filtered historical simulation (per-asset EWMA rescaling), GARCH(1,1) with Student-$t$
innovations (1000-day window, refitted monthly), GARCH-EVT \citep{mcneil2000} with a generalised Pareto tail
above the 10\% threshold of the standardised residuals, and DCC-GARCH-$t$ (univariate GARCH refitted monthly,
DCC parameters yearly). Multivariate models produce joint scenarios, so they revalue the non-linear book too.

% ------------------------------------------------------------------------------------------------
\section{Data-only baseline}
\label{sec:met-dataonly}

To measure what the physics contributes (RQ3), the $d$-dimensional network is trained without the PIDE. It is
identical to the PINN in architecture, optimiser, learning-rate schedule, number of steps and normaliser
estimator, and it is trained by maximum likelihood on $N$ exact samples $(x_i, t_i)$:
\begin{equation}
  \mathcal{L}_{\mathrm{NLL}} = -\frac1N \sum_i \Big[u(x_i, t_i) - \log Z(t_i)\Big], \qquad
  \log Z(t) = \log \mathbb{E}_{q}\!\Big[e^{\,u(X,t) - \log q(X \mid t)}\Big],
\end{equation}
where $\log Z$ is estimated on fresh proposal draws at every step, exactly as the PINN's normalisation term.
The samples are spread evenly over 16 times: 13 log-spaced in $[t_0, T]$ plus the three evaluation horizons.
20\% are held out, and the state with the best validation likelihood (on fixed proposal draws) is kept. $N$
ranges over $10^3$, $10^4$, $10^5$ and $10^6$, with three seeds each. Monte Carlo on the same samples (those at
each horizon) is reported as the purely empirical alternative.

% ------------------------------------------------------------------------------------------------
\section{Evaluation protocol}
\label{sec:met-eval}

\paragraph{Portfolios.} For each dimension: the equal-weight portfolio, 15 long-only portfolios with
Dirichlet(1) weights and 4 long--short portfolios (weights summing to one), with a fixed seed. Errors are
reported as the median over these 20 portfolios.

\paragraph{Accuracy.} The headline metric is the relative error of $\ES_{99}$ at 1, 5 and 10 days against the
exact value; VaR and the 97.5\% level are also reported. Density accuracy is split into bulk and tail, which
answers RQ1 directly. In 1-D the bulk error is the relative $L^1$ error between the 1\% and 99\% quantiles and the
tail error the mean absolute error of $\log p$ beyond them. In $d$ dimensions both are evaluated on exact samples:
the mean absolute error of $\log p$ over all samples, and over the samples in the 1\% loss tail of the equal-weight
portfolio.

\paragraph{Sampling controls.} Importance sampling adds sampling error to the network's error. Two controls are
computed on the same draws: the same estimator with the \emph{exact} density (the IS floor, the best any density
could do with this proposal and budget), and Monte Carlo with the same number of draws. All methods in a
comparison use common random numbers, so differences between them are model differences.

\paragraph{Backtests.} For each forecast series, with a breach on day $t$ when the realised P\&L is below
$-\VaR_t$:
\begin{itemize}
  \item \textbf{Kupiec} unconditional coverage: likelihood-ratio test that the breach rate is $1 - \alpha$.
  \item \textbf{Christoffersen} conditional coverage: Kupiec plus a test that breaches are independent from one
        day to the next (a first-order Markov alternative).
  \item \textbf{FZ0} score of \citet{patton2019}, a strictly consistent scoring function for the pair
        (VaR, ES): with $v_t = -\VaR_t$, $e_t = -\ES_t$ and $p = 1-\alpha$,
        \begin{equation}
          S(y_t; v_t, e_t) = -\frac{\mathbf{1}\{y_t \le v_t\}(v_t - y_t)}{p\, e_t} + \frac{v_t}{e_t} + \log(-e_t) - 1,
        \end{equation}
        averaged over the test period (lower is better). It is the only one of these criteria that evaluates ES.
  \item \textbf{Diebold--Mariano} test of equal mean FZ0 loss between each benchmark and the PINN, with a
        Newey--West variance.
\end{itemize}

\paragraph{Failure cases.} Each failure experiment changes one setting of the $d = 2$ model and keeps all
others: a crisis jump intensity $\lambda = 50$/yr with the default proposal ($K = 3$, which leaves 14\% of the
10-day mass in states with more jumps) and with $K = 10$, and a start time ten times closer to the singular
initial condition ($t_0 = 0.05$ day). A far-tail probe evaluates $\log p - \log q$ for the PINN and the exact
density along 2000 random directions up to 150 whitened units.

\paragraph{Reproducibility.} Every experiment is a script driven by a YAML configuration. Each run writes its
figures, its metrics and the code version, library versions and device that produced them. The notebooks only
call these scripts, and the test suite checks the exact density, the reference solvers, the PINN residual
(which vanishes for the exact density as the quadrature is refined) and the data-only likelihood.
